\documentclass{situationist-preprint}

\title{Dreams as Conscious Experiences}
\msid{MND-25-a}
\msfamily{MND-25}
\msdiscipline{Philosophy of Mind}
\mstype{Principal manuscript}
\mskeywords{dreaming; sleep; dream reports; Malcolm; Dennett; lucid dreaming; REM sleep; memory reconstruction}

\begin{document}
\maketitle

\begin{abstract}
Are dreams experiences? Malcolm argued that the claim that one had experiences while asleep is unverifiable and that dreams are nothing more than the dispositions to tell dream reports on waking. Dennett proposed that dream memories might be composed at the moment of waking, with no experience during sleep. We argue that dreams are conscious experiences during sleep, though dream reports are partly reconstructed on waking. The argument rests on evidence that ties the content of dreams to events during sleep: the correspondence between the reported duration and order of dream events and the timing of sleep, the eye signals and real-time answers given by lucid dreamers from REM sleep, the enactment of dream content in REM sleep behavior disorder, and the prediction of dream reports from regional brain activity before waking. A composition-on-waking theory cannot explain this evidence; a theory on which experiences occur during sleep and are imperfectly remembered can. The reconstructive element in reports explains their gaps and their narrative smoothing without denying the experiences.
\end{abstract}

\section{The problem}

On waking, people often report having seen, heard and done things while asleep. The common view, and that of most dream science, is that these reports are memories of experiences had during sleep. \citet{Malcolm1959} argued that this view has no content: since a sleeping person cannot judge or report anything, nothing could verify that she had experiences, and \enquote{she dreamt that $p$} means only that on waking she is inclined to tell that $p$ as a dream. \citet{Dennett1976} gave the sceptical position an empirical form: on the \enquote{cassette theory}, dream contents are stored unconsciously during sleep and composed into a narrative on waking, without ever being experienced, and the received view and the cassette theory could be distinguished only by evidence about timing and about the states of the dreamer during sleep. A different line, defended by \citet{Sosa2005} and \citet{Ichikawa2009}, holds that dreams are imaginings rather than illusory perceptions and beliefs; this view is compatible with dreams being experiences and we do not need to resolve it. \citet{Windt2015} defends the view that dreaming is a form of experience, conceived as immersive spatiotemporal hallucination during sleep, and assesses the reliability of dream reports.

We argue that dreams are conscious experiences during sleep, and that dream reports are partly reconstructed on waking.

\section{Framework}

\begin{definition}[Dream; dream report]\label{d:dream}
A \emph{dream report} is a report, made on waking, of apparent experiences during sleep. A \emph{dream}, on the experiential hypothesis, is a sequence of conscious experiences during sleep of which dream reports are, when accurate, reports.
\end{definition}

\begin{definition}[Composition hypothesis]\label{d:comp}
The \emph{composition hypothesis} is that dream reports are formed from information stored without experience during sleep and composed into apparent memories at or after waking.
\end{definition}

\begin{principle}[Inference to the best explanation]\label{p:ibe}
Of two hypotheses about the occurrence of experiences, the one that better explains the correlations between reports and independently measured events is to be preferred.
\end{principle}

\begin{principle}[Timing]\label{p:time}
If a reported conscious event co-varies in time with a measured event during sleep, and the subject can, from within the reported episode, produce a response that is measured at that time and that answers to its content, the reported event occurred at that time.
\end{principle}

\section{The argument}

\begin{lemma}\label{l:evidence}
The following correlations hold. (i) Reports from awakenings in REM sleep are more frequent and more vivid than from other stages, and the reported duration of dream episodes grows with the time spent in REM sleep before waking \citep{DementKleitman1957}. (ii) Lucid dreamers can signal, by agreed eye movements recorded in REM sleep, that they have become lucid, and later report having made the signals in the dream \citep{LaBerge1981}; they can answer questions put to them during REM sleep, such as simple arithmetic problems, by eye movements or facial muscle signals, and report on waking having heard the question in the dream \citep{Konkoly2021}. (iii) Patients with REM sleep behavior disorder, who lack the normal paralysis of REM sleep, act out movements that correspond to the content they report on waking \citep{Schenck1986}. (iv) Regional activity in a posterior cortical zone before awakening predicts whether a subject will report having dreamed and aspects of what she will report \citep{Siclari2017}.
\end{lemma}

The lemma summarizes the evidence; the cited studies give the details and their limits.

\begin{theorem}\label{t:main}
Dreams are conscious experiences during sleep.
\end{theorem}

\begin{proof}
By Principle~\ref{p:time} and Lemma~\ref{l:evidence}(ii), the experiences reported by lucid dreamers as including the hearing of a question and the making of a signal occurred during sleep, at the times at which the signals were recorded, since the signals answer to the content of the questions and were made before waking. So at least some dreams are experiences during sleep, contrary to the composition hypothesis (Definition~\ref{d:comp}). The composition hypothesis must explain (i), (iii) and (iv) by unconscious processes that mimic the timing and content of experience, and must treat the lucid cases as exceptions; the experiential hypothesis explains all four by the same thing, experiences during sleep that are later reported. Lucid dreams differ from non-lucid ones in the presence of insight that one is dreaming, not in the presence of imagery, scenes and agency, and the measured correlations of (i), (iii) and (iv) hold for non-lucid dreams. By Principle~\ref{p:ibe}, dreams in general are conscious experiences during sleep.
\end{proof}

\begin{proposition}\label{p:recon}
Dream reports are partly reconstructed on waking.
\end{proposition}

\begin{proof}
Dream recall decays within minutes of waking unless rehearsed, and subjects who report on the same dream at successive times give reports that differ, with later reports more coherent and narratively ordered. Reports also include content from the waking context and from expectation. These are the marks of memory reconstruction, which supplies schema-driven connections to fill gaps in encoded material. Since dream experiences are encoded during a state in which memory consolidation of episodic material is weak, reconstruction is to be expected in their reports. So reports are partly reconstructed.
\end{proof}

Proposition~\ref{p:recon} explains why the question whether one dreams in color, or how long a dream lasted, is unreliably answered, and is the ground on which MND-20 classes such reports among judgments about non-current states, where introspective reliability does not extend.

\section{Objections and replies}

\begin{objection}[Malcolm]\label{o:malc}
The claim that a sleeper has experiences cannot be verified, since sleep excludes the judgment and report that verification would require \citep{Malcolm1959}.
\end{objection}

\begin{reply}
The objection assumes that sleep excludes report. Lemma~\ref{l:evidence}(ii) shows that it does not: lucid dreamers report from within sleep. And the claim is verifiable by the indirect methods by which we verify other claims about unobservable states, which Principle~\ref{p:ibe} records.
\end{reply}

\begin{objection}[Lucidity]\label{o:lucid}
Lucid dreams may be a mixed state, partly waking; the conclusion should not be extended to ordinary dreams.
\end{objection}

\begin{reply}
Lucid signals are recorded in polysomnographically verified REM sleep. And the extension to non-lucid dreams does not rest on the lucid cases alone, but on the correlations of Lemma~\ref{l:evidence}(i), (iii) and (iv), which concern ordinary dreams.
\end{reply}

\begin{objection}[Imagination]\label{o:imag}
Dreams are imaginings, not perceptual experiences; so they are not experiences in the sense intended \citep{Ichikawa2009}.
\end{objection}

\begin{reply}
Imaginings are conscious states with phenomenal character. Whether dreams are better classed with imaginings or with hallucinations bears on the kind of experience they are, not on whether they are experiences.
\end{reply}

\section{What the result does not establish}

The argument does not show that every dream report corresponds to a dream, nor that dreams occur only in REM sleep; reports from other stages are common and may be partly reconstructed from less vivid experience. It does not show that dreams are imaginings rather than hallucinatory experiences. But it bears on the kind of experience dreams are. Dreams are imagistic experiences generated without the perceived objects; MND-16 holds that hallucinations are experiences of a different kind from veridical perceptions, and on the present account hallucinations are of the same broad kind as dreams. MND-23 notes that some mammals and birds show sleep stages resembling human REM sleep; the present result does not by itself show that they dream.

\section{Conclusion}

The correlations between dream reports and events during sleep, and above all the reports given from within sleep by lucid dreamers, show that dreams are conscious experiences during sleep. The reports by which we know them are partly reconstructed on waking, which explains their gaps and their unreliability on matters of detail, without impugning the experiences they report.

\begin{thebibliography}{99}\small
\bibitem[Dement and Kleitman(1957)]{DementKleitman1957} Dement, W. and Kleitman, N. (1957). The relation of eye movements during sleep to dream activity: An objective method for the study of dreaming. \emph{Journal of Experimental Psychology}, 53(5), 339--346.
\bibitem[Dennett(1976)]{Dennett1976} Dennett, D.~C. (1976). Are dreams experiences? \emph{The Philosophical Review}, 85(2), 151--171.
\bibitem[Ichikawa(2009)]{Ichikawa2009} Ichikawa, J. (2009). Dreaming and imagination. \emph{Mind \& Language}, 24(1), 103--121.
\bibitem[Konkoly et~al.(2021)]{Konkoly2021} Konkoly, K.~R., Appel, K., Chabani, E., et~al. (2021). Real-time dialogue between experimenters and dreamers during REM sleep. \emph{Current Biology}, 31(7), 1417--1427.
\bibitem[LaBerge et~al.(1981)]{LaBerge1981} LaBerge, S.~P., Nagel, L.~E., Dement, W.~C. and Zarcone, V.~P. (1981). Lucid dreaming verified by volitional communication during REM sleep. \emph{Perceptual and Motor Skills}, 52(3), 727--732.
\bibitem[Malcolm(1959)]{Malcolm1959} Malcolm, N. (1959). \emph{Dreaming}. Routledge \& Kegan Paul.
\bibitem[Schenck et~al.(1986)]{Schenck1986} Schenck, C.~H., Bundlie, S.~R., Ettinger, M.~G. and Mahowald, M.~W. (1986). Chronic behavioral disorders of human REM sleep: A new category of parasomnia. \emph{Sleep}, 9(2), 293--308.
\bibitem[Siclari et~al.(2017)]{Siclari2017} Siclari, F., Baird, B., Perogamvros, L., et~al. (2017). The neural correlates of dreaming. \emph{Nature Neuroscience}, 20(6), 872--878.
\bibitem[Sosa(2005)]{Sosa2005} Sosa, E. (2005). Dreams and philosophy. \emph{Proceedings and Addresses of the American Philosophical Association}, 79(2), 7--18.
\bibitem[Windt(2015)]{Windt2015} Windt, J.~M. (2015). \emph{Dreaming: A Conceptual Framework for Philosophy of Mind and Empirical Research}. MIT Press.
\end{thebibliography}
\end{document}
